\documentclass[11pt]{article}
\usepackage[margin=1in]{geometry}
\usepackage{booktabs}
\usepackage{amsmath}
\usepackage{hyperref}
\title{Active SolvAI: learning what to measure in alchemical hydration calculations}
\author{Gal Oren \and Michael Levitt}
\date{}
\begin{document}
\maketitle

\begin{abstract}
Accurate hydration free energies require target-specific sampling, whereas structure-only models can fail under chemical shift. Active SolvAI asks whether sparse response observations from the query molecule can update a transferable SolvAI prior and recover either the same-Hamiltonian alchemical response or the experimental endpoint at lower measured cost. This manuscript is a living preregistered report: quantitative claims are inserted only after their generating analysis has been frozen and completed.
\end{abstract}

\section{Introduction}
SolvAI established that structure-predicted solvent responses add information beyond matched molecular features and experimental labels. It did not establish that high-fidelity response can be recovered from structure alone, nor that uncertainty can guide a simulation. We therefore test a different proposition: actual sparse observations may reveal molecule-specific departures from the structure prior.

\section{Results}
\subsection{Actual-observation gate}
\textbf{PHASE1\_RESULT\_PENDING.} The primary estimand is the paired held-out endpoint-error change obtained from actual-minus-prior response residuals relative to frozen SolvAI. Same-Hamiltonian reconstruction and experimental prediction are reported separately.

\subsection{Retrospective replay}
\textbf{REPLAY\_RESULT\_PENDING.} This section will be populated only if the preregistered Phase 1 go criterion is met.

\subsection{Prospective sentinels}
\textbf{PILOT\_RESULT\_PENDING.} This section will be populated only after a replay policy and its comparators are frozen.

\section{Discussion}
The study is designed to support a positive, neutral or negative conclusion without changing its scientific question. Any benefit for same-Hamiltonian reconstruction will not by itself be described as improved agreement with experiment; any endpoint benefit without curve reconstruction will be described as an empirical residual diagnostic rather than successful quadrature.

\section{Methods}
The parent SolvAI artifact, cohorts, splits and response campaigns are versioned by SHA-256. All preprocessing, calibration and model selection occur inside molecule-held-out training data. Destructive controls preserve response distributions while shuffling molecule alignment. Every attempted run, including failures, is included in the compute ledger.

\end{document}

